%Time-stamp: <qua nov 12 11:53:52 2025 yair@asterix>
%\usepackage{bbm}%\1
\usepackage{dsfont} % para \mathds{1}
\def\mathbbm#1{\mathds{#1}}

\usepackage{amssymb,amsthm,mathtools}
\mathtoolsset{
  showonlyrefs=true % or false in draft mode
}
\usepackage{xparse}

\def\To{\Rightarrow}

%------------------------------------------------------------------%
%    CONJUNTOS
%%-------------------------------------------------------------------
\DeclareRobustCommand{\C}[0]{\mathbb{C}}
\DeclareRobustCommand{\F}[0]{\mathbb{F}}
\DeclareRobustCommand{\N}[0]{\mathbb{N}}
\DeclareRobustCommand{\Q}[0]{\mathbb{Q}}
\DeclareRobustCommand{\R}[0]{\mathbb{R}}
\DeclareRobustCommand{\Z}[0]{\mathbb{Z}}
\DeclareRobustCommand{\P}[0]{\mathbb{P}}
\DeclareRobustCommand{\E}[0]{\mathbb{E}}
%%---------------------------------------------------------------------
%   DELIMITADORES
%Now if you use the starred version of the macro you get the desired
% result: \setc*{x \in X}{x > \frac{1}{2}} You can even adjust to a
% different size if you think that the braces are too big, with the
% syntax \setc[\Big]{x \in X}{x > \frac{1}{2}} 
%%---------------------------------------------------------------------
\DeclarePairedDelimiter\card{\lvert}{\rvert}
\DeclarePairedDelimiter\paren{\lparen}{\rparen}
\DeclarePairedDelimiter\colch{\lbrack}{\rbrack}
\DeclarePairedDelimiter\teto{\lceil}{\rceil}
\DeclarePairedDelimiter\chao{\lfloor}{\rfloor}
\DeclarePairedDelimiter\norma{\lVert}{\rVert}

\DeclareRobustCommand{\Card}[1]{\card*{\,#1\,} }
\DeclareRobustCommand{\Norma}[1]{\norma*{\,#1\,} }

%conjunto
\DeclarePairedDelimiterX\conj[1]{\lbrace}{\rbrace}{\,#1\,}
\DeclarePairedDelimiterX\chaves[1]{\lbrace}{\rbrace}{\,#1\,}

% condicionais
\DeclarePairedDelimiterX\conjc[2]{\lbrace}{\rbrace}%
% {#1 \;\delimsize\mid\; #2}
{\,#1 \;\delimsize\vert\; #2\,}
%{\,#1 \middle|  #2\,}
\DeclarePairedDelimiterX\chavesc[2]{\lbrace}{\rbrace}%
% {#1 \;\delimsize\mid\; #2}
{\,#1 \;\delimsize\vert\; #2\,}
%{\,#1 \middle|  #2\,}
\DeclarePairedDelimiterX\parenc[2]{\lparen}{\rparen}%
%{\,#1 \middle|  #2\,}
{\,#1 \;\delimsize\vert\; #2\,}
%{#1 \;\delimsize\mid\; #2}
\DeclarePairedDelimiterX\colchc[2]{\lbrack}{\rbrack}%
%{\,#1 \middle|  #2\,}
{\,#1 \;\delimsize\vert\; #2\,}
%{#1 \;\delimsize\mid\; #2}
%---------------------------------------------------------------------
%---------------------------------------------------------------------
%   PROBABILIDADE
%---------------------------------------------------------------------
\DeclareRobustCommand{\prob}[2][{\protect{{\P}}}]%
{\ensuremath{\mathbb{#1} \paren{\,#2\,}}}
%\DeclareRobustCommand{\probE}[2][\protect\operatorname{\P}]%
\DeclareRobustCommand{\probE}[2][\protect{\P}]%
{\ensuremath{\mathbb{#1} \colch{\,#2\,}}}
\DeclareRobustCommand{\pLei}[3][\protect\operatorname{\P}]%
{\ensuremath%
  \protect\operatorname{{\mathbb {#1}}_{#2}}%
  \paren{#3}}%
\DeclareRobustCommand{\pLeiE}[3][\protect\operatorname{\P}]%
{\ensuremath%
  \protect\operatorname{{\mathbb {#1}}_{#2}} \colch{#3}}%
% condicional -----------------------------------
\DeclareRobustCommand{\probC}[3][\protect\operatorname{\P}]%
{\ensuremath{\mathbb {#1} \parenc{\,#2\,}{\,#3\,} }}
% condicional descrição de evento
\DeclareRobustCommand{\probCE}[3][\protect\operatorname{\P}]%
{\ensuremath{\mathbb{#1\,} \colchc{\,#2\,}{\,#3\,} }}
\DeclareRobustCommand{\PX}[2][\protect\operatorname{\P}]%
{\ensuremath%
  \protect\operatorname{\mathbb {#1}}%
  \displaylimits_{#2}}
% P_{}()-----------------------------------
\DeclareRobustCommand{\probX}[3][\protect\operatorname{\P}]%
{\ensuremath%
  \protect\operatorname{\mathbb {#1}}%
  \displaylimits_{#2} \paren{#3}}
% P_{}[]-----------------------------------
\DeclareRobustCommand{\probXE}[3][\protect\operatorname{\P}]%
{\ensuremath%
  \protect\operatorname{\mathbb {#1}}%
  \displaylimits_{#2} \colch{#3}}
% P_{}( | )-----------------------------------
\DeclareRobustCommand{\probCX}[4][\protect\operatorname{\P}]%
{\ensuremath%
  \protect\operatorname{\mathbb {#1}}%
  \displaylimits_{#2} \parenc{#3}{#4} }
% P_{}[ | ]-----------------------------------
\DeclareRobustCommand{\probCXE}[4][\protect\operatorname{\P}]%
{\ensuremath%
  \protect\operatorname{\mathbb {#1}}%
  \displaylimits_{#2} \colchc{#3}{#4} }
%
% probabilidade com ajuste de delimitadores
%
\DeclareRobustCommand{\Prob}[2][\protect\operatorname{\P}]%
{\ensuremath{\mathbb{#1} \paren*{\,#2\,}}}
\DeclareRobustCommand{\ProbE}[2][\protect\operatorname{\P}]%
{\ensuremath{\mathbb{#1} \colch*{\,#2\,}}}
% condicional -----------------------------------
\DeclareRobustCommand{\ProbC}[3][\protect\operatorname{\P}]%
{\ensuremath{\mathbb {#1} \parenc* {\,#2\,}{\,#3\,} }}
% condicional descrição de evento
\DeclareRobustCommand{\ProbCE}[3][\protect\operatorname{\P}]%
{\ensuremath{\mathbb{#1\,} \colchc* {\,#2\,}{\,#3\,} }}
\DeclareRobustCommand{\PX}[2][\protect\operatorname{\P}]%
{\ensuremath%
  \protect\operatorname{\mathbb {#1}}%
  \displaylimits_{#2}}
% P_{}()-----------------------------------
\DeclareRobustCommand{\PLei}[3][\protect\operatorname{\P}]%
{\ensuremath%
  \protect\operatorname{{\mathbb {#1}}_{#2}}%
  \paren*{#3}}%
\DeclareRobustCommand{\PLeiE}[3][\protect\operatorname{\P}]%
{\ensuremath%
  \protect\operatorname{{\mathbb {#1}}_{#2}} \colch*{#3}}%
\DeclareRobustCommand{\ProbX}[3][\protect\operatorname{\P}]%
{\ensuremath%
  \protect\operatorname{\mathbb {#1}}%
  \displaylimits_{#2} \paren*{#3}}
% P_{}[]-----------------------------------
\DeclareRobustCommand{\ProbXE}[3][\protect\operatorname{\P}]%
{\ensuremath%
  \protect\operatorname{\mathbb {#1}}%
  \displaylimits_{#2} \colch*{#3}}
% P_{}( | )-----------------------------------
\DeclareRobustCommand{\ProbCX}[4][\protect\operatorname{\P}]%
{\ensuremath%
  \protect\operatorname{\mathbb {#1}}%
  \displaylimits_{#2} \parenc* {#3}{#4} }
% P_{}[ | ]-----------------------------------
\DeclareRobustCommand{\ProbCXE}[4][\protect\operatorname{\P}]%
{\ensuremath%
  \protect\operatorname{\mathbb {#1}}%
  \displaylimits_{#2} \colchc* {#3}{#4} }
%%---------------------------------------------------------------------
%%---------------------------------------------------------------------
%% indicador de evento-----------------------------------
\DeclareRobustCommand{\1}[1]{\ensuremath \mathbbm{1}_{#1}}
%
%% esperança-----------------------------------
\DeclareRobustCommand{\Es}[2][\protect\operatorname{\mathbb{E}}]%
{\ensuremath%
  \protect\operatorname{\mathbb {#1}}{\,#2\,} }
\DeclareRobustCommand{\EX}[3][\protect\operatorname{\mathbb{E}}]%
{\ensuremath%
  \protect\operatorname{\mathbb {#1}_{#2}}{\,#3\,} }
%% esperança expressão-----------------------------------
\DeclareRobustCommand{\eE}[2][\protect\operatorname{\mathbb{E}}]%
{\ensuremath%
  {#1\,}\colch{\,#2\,} }
%% esperança ondicional
\DeclareRobustCommand{\eC}[3][\protect\operatorname{\mathbb{E}}]%
{\ensuremath%
  {#1\,}\colchc {#2}{#3} }
% variancia
\DeclareRobustCommand{\var}[2][\protect\operatorname{\mathrm{Var}}]{%
  \ensuremath {#1\,} \colch{\,#2\,}}
%
% esperança com ajuste de delimitadores
%% esperança expressão-----------------------------------
\DeclareRobustCommand{\EE}[2][\protect\operatorname{\mathbb{E}}]%
{\ensuremath%
  {#1\,}\colch*{\,#2\,} }
%% esperança ondicional
\DeclareRobustCommand{\EC}[3][\protect\operatorname{\mathbb{E}}]%
{\ensuremath%
  {#1\,}\colchc* {#2}{#3} }
% variancia
\DeclareRobustCommand{\Var}[2][\protect\operatorname{\mathrm{Var}}]{%
  \ensuremath {#1\,} \colch*{\,#2\,}}

%%--------------------------------------------------------------
%%-ALIAS-----------------------------------------------------
% \NewDocumentCommand\mycommand{}{somecodehere}
% \NewDocumentCommand\varpower{mO{n}}{#1^{#2}}
%
\NewDocumentCommand\eps{}{\ensuremath \varepsilon}
\NewDocumentCommand\rest{}{\ensuremath \!\!\restriction}
\NewDocumentCommand\qeg{}{\hfill \ensuremath{\lozenge}}
% 
% \RenewDocumentCommand\mtctitle{}{Conteúdo}
%
%%%% %conjunto vazio% %%%%%%%%%%%%%%%
\RenewDocumentCommand\emptyset{}{\ensuremath \varnothing}
% 
% %em Distribuições%%%
\NewDocumentCommand\indist{m}{\mathop{\in_{\mathtt{#1}}}}
%\NewDocumentCommand\inuni{}{\mathop{\in_{\mathcal{U}}}}
\NewDocumentCommand\inuni{}{\mathop{\in_{\mathrm{R}}}}
%\NewDocumentCommand\getuni{}{\xleftarrow{\mathrm{R}}}
\NewDocumentCommand\getuni{}{\xleftarrow{\mathrm{{}_R}}}
\NewDocumentCommand\inber{m}{\mathop{\in_{\tiny{\mathrm{b}_{#1}}}}}
\NewDocumentCommand\inbin{m}{\mathop{\in_{\tiny{\mathrm{b}_{#1}}}}}
\NewDocumentCommand\ingeo{m}{\mathop{\in_{{G}_{#1}}}}
% 
% 
%\NewDocumentCommand\eqdef{}{\coloneqq}
\NewDocumentCommand\partes{m}{%
  {\raisebox{.1\baselineskip}{\large\ensuremath{\wp}}}(#1)}
\RenewDocumentCommand\:{}{\colon}
\NewDocumentCommand\<{}{\leqslant}
\RenewDocumentCommand\>{}{\geqslant}
%no. de Stirling
\DeclareRobustCommand{\stirling}{\genfrac\{\}{0pt}{}}
%--------------------------------------------------------------------
\hyphenation{}
% --------------------------------------------------------------------

\makeatletter
\newcommand*{\dotgrande}{}% Check if undefined
\DeclareRobustCommand*{\dotgrande}{%
  \mathbin{\mathpalette\dotgrande@{}}%
}
\newcommand*{\dotgrande@scalefactor}{.5}
\newcommand*{\dotgrande@widthfactor}{1.15}
\newcommand*{\dotgrande@}[2]{%
  % #1: math style
  % #2: unused
  \sbox0{$#1\vcenter{}$}% math axis
  \sbox2{$#1\cdot\m@th$}%
  \hbox to \dotgrande@widthfactor\wd2{%
    \hfil
    \raise\ht0\hbox{%
      \scalebox{\dotgrande@scalefactor}{%
        \lower\ht0\hbox{$#1\bullet\m@th$}%
      }%
    }%
    \hfil
  }%
}
\makeatother



%%% Local Variables:
%%% mode: latex
%%% TeX-master: t
%%% End:
